\newpage
\phantomsection
\label{prf:cousin}
\LRAProofFor{lem:cousin}
\begin{remark*}[Return]\hyperref[lem:cousin]{Return to Lemma}\end{remark*}
\begin{lemma*}[Cousin's Lemma]
For every gauge \(\delta\) on \([a,b]\), there exists a \(\delta\)-fine tagged
partition of \([a,b]\).
\end{lemma*}
\begin{proof}[Professional Standard Proof]
\LRAProofBodyStart
Assume no \(\delta\)-fine tagged partition exists. Bisect \([a,b]\). At least
one half must also admit no \(\delta\)-fine tagged partition, because
partitions of both halves could be concatenated. Iterating gives nested closed
intervals \([a_n,b_n]\) with lengths tending to \(0\), each admitting no
\(\delta\)-fine tagged partition. By the nested interval property their
intersection is a single point \(c\). Since \(\delta(c)>0\), some interval
\([a_N,b_N]\) is contained in
\((c-\delta(c),c+\delta(c))\). The one-slab partition of \([a_N,b_N]\), tagged
at \(c\), is then \(\delta\)-fine, a contradiction.
\end{proof}
\begin{proof}[Detailed Learning Proof]
\LRAProofBodyStart
The only possible way a gauge could fail is if every attempted chopping has a
bad subinterval. Repeated bisection traps those bad subintervals around one
point. But the gauge is positive at that point, so sufficiently small trapped
intervals fit inside its gauge ball and are valid after all. This contradiction
proves non-vacuity.
\end{proof}
\begin{remark*}[Proof structure]
Assume failure, repeatedly bisect to produce nested bad intervals, then use
compactness and positivity of the gauge at the limiting point to build the
forbidden one-slab partition.
\end{remark*}
\begin{dependencies}
\begin{itemize}
  \item \hyperref[def:gauge-on-interval]{Gauge on an Interval}.
  \item \hyperref[def:delta-fine-tagged-partition]{Delta-Fine Tagged Partition}.
\end{itemize}
\end{dependencies}
\clearpage
